\subsection{Sprint 2}

A \textit{sprint} 2 foi a mais curta do semestre, com 16 dias, de 14 a 30 de
setembro. Planejamos três épicos: Cadastro de Usuário, Participação de Eventos
e Notificações. Cada um ficou com um dos três \textit{squads}, e o meu
assumiu Participação de Eventos. Reforcei com os colegas que a divisão servia
para aproximar grupos menores e facilitar o acompanhamento. Perto da entrega,
todos deveriam ajudar nas pendências, mesmo que fossem de outro
\textit{squad}.

Também assumi a elaboração do plano de comunicação, o último artefato de
gestão que faltava, para não deixá-lo pendente até a entrega final. Durante
a \textit{sprint}, concentrei minha
atuação no acompanhamento dos prazos, das \textit{dailies}, das presenças e
dos atrasos, além da preparação desse plano.

Entregamos todas as tarefas previstas, e o produto agradou aos clientes.
Meu \textit{squad} concluiu o fluxo de participação em eventos e ajudou nas
últimas pendências dos demais grupos. Como o prazo era curto, tivemos de
definir datas mais rígidas para a conclusão das tarefas perto da entrega e
pedir um esforço conjunto da equipe.

As dependências entre tarefas exigiram mais atenção, principalmente quando
envolviam \textit{squads} diferentes. Uma atividade podia estar pronta para
começar e ainda depender de um componente ou de uma alteração de outra
equipe. Levantamos essas relações informalmente e conversamos com os
integrantes de AGES III e IV de cada grupo para priorizar os \acp{pr}
bloqueantes e orientar a ordem de implementação.

Decidimos documentar as dependências no início de cada \textit{sprint}.
A proposta era listar as tarefas e indicar quais atividades cada uma
bloqueava, distinguindo dependências internas e externas ao \textit{squad}.
Queríamos considerar essas relações ao definir as prioridades de cada grupo.

Para a etapa seguinte, restavam mais três épicos, novamente distribuídos
entre os \textit{squads}. Já sabíamos que feriados e a semana acadêmica
reduziriam o tempo disponível. Planejamos acompanhar os prazos com mais
atenção, registrar as dependências e usar os marcadores de prioridade para
orientar a equipe.
